% =====================================================================
% Appendix: hardware, decoding and prompt protocol.
% Every value here is taken from the pinned run configs and the vLLM
% evaluation contract, not reconstructed from memory.
%   configs/local_models.json, configs/gpt_oss_models.json,
%   evaluation/vllm/README.md, finexam_eval/data.py::prompt_for
% =====================================================================

\subsection{Hardware Configuration}
\label{app:hardware}

Hardware assignment for each locally deployed model followed its memory footprint under
bfloat16 with a $32{,}768$-token context. The five finance-specialized models and the two
open-weight general models were served with vLLM on NVIDIA H100 accelerators.
DianJin-R1-32B is the only model in the suite that does not fit on one card: at $32$B
parameters it needs tensor parallelism of two and was run on $2\times$H100 with $160$~GB of
host memory at batch size two. Hawkish-8B, Fin-O1-14B, ODA-Fin-RL-8B and Fin-R1-7B each ran
on a single H100 with $64$~GB of host memory at tensor parallelism one and batch size four.

The two GPT-OSS mixture-of-experts models were also served on a single H100 each, using their
native MXFP4 weights with a bfloat16 compute dtype. GPT-OSS-20B ran at batch size four with a
GPU memory utilisation target of $0.90$, and GPT-OSS-120B at batch size one with $0.95$, the
larger reservation being required because the $117$B expert set leaves little headroom for the
key-value cache on one card. Because MXFP4 and the associated Triton kernels are not available
in older releases, the GPT-OSS runs used the same pinned environment as the rest of the suite
rather than a legacy stack, which is what makes a single shared software configuration possible
across all seven models.

The remaining ten systems are proprietary or API-only and were evaluated through their official
endpoints, which removes any local GPU requirement. DeepSeek-R1 belongs to this group despite
its MIT licence, because a $671$B mixture of experts cannot be served on the hardware described
above.

Table~\ref{tab:local-config} lists the pinned revision and serving configuration for every
locally deployed model. Pinning revisions matters here because several of these checkpoints have
been updated since evaluation, and results reported against a moving \texttt{main} branch are
not reproducible.

\begin{table*}[t]
\centering\small
\setlength{\tabcolsep}{4pt}
\begin{tabular}{llllrrr}
\toprule
\textbf{Model} & \textbf{Hugging Face identifier} & \textbf{Revision} & \textbf{Precision}
& \textbf{TP} & \textbf{GPUs} & \textbf{Batch} \\
\midrule
GPT-OSS-120B   & \texttt{openai/gpt-oss-120b}          & \texttt{b5c939de} & MXFP4 / BF16 & 1 & 1 & 1 \\
GPT-OSS-20B    & \texttt{openai/gpt-oss-20b}           & \texttt{6cee5e81} & MXFP4 / BF16 & 1 & 1 & 4 \\
\midrule
DianJin-R1-32B & \texttt{DianJin/DianJin-R1-32B}       & \texttt{f2976414} & BF16 & 2 & 2 & 2 \\
Fin-O1-14B     & \texttt{TheFinAI/Fin-o1-14B}          & \texttt{dadd79b8} & BF16 & 1 & 1 & 4 \\
ODA-Fin-RL-8B  & \texttt{OpenDataArena/ODA-Fin-RL-8B}  & \texttt{948c22ea} & BF16 & 1 & 1 & 4 \\
Hawkish-8B     & \texttt{mukaj/Llama-3.1-Hawkish-8B}   & \texttt{bd4968f5} & BF16 & 1 & 1 & 4 \\
Fin-R1-7B      & \texttt{SUFE-AIFLM-Lab/Fin-R1}        & \texttt{1ae9364b} & BF16 & 1 & 1 & 4 \\
\bottomrule
\end{tabular}
\caption{Serving configuration for the seven locally deployed models. All use a
$32{,}768$-token context window and an $8{,}192$-token completion budget on NVIDIA H100
accelerators. Revisions are truncated to eight characters. Every model was pinned to an
immutable commit rather than to a branch.}
\label{tab:local-config}
\end{table*}

\paragraph{Software stack.}
All local evaluations ran in one isolated environment on Python 3.11 with vLLM 0.10.0,
Transformers 4.53.2, Hugging Face Hub 0.33.4, tokenizers 0.21.2 and safetensors 0.5.3.
Inference was offline, with inherited Hugging Face credentials cleared at process start. Each
run records an execution hash covering the backend, the model and tokenizer revisions, the full
prompt and sampling configuration, the dtype, the context length, the batch size, the tensor
parallel degree, the sampling seed and the installed vLLM, Transformers, Torch and tokenizers
versions. Resume is refused when that hash does not match, so a run cannot be silently completed
under a changed environment.

\paragraph{Decoding.}
Every model, local or Proprietary, used a single unified protocol: temperature $0$, top-$p$ $1$,
seed $202607$, and a completion budget of $8{,}192$ tokens. We state explicitly that this is a
protocol choice rather than a reproduction of vendor defaults. Fin-O1-14B and ODA-Fin-RL-8B ship
official generation configs of temperature $0.6$, top-$p$ $0.95$ and top-$k$ $20$, and Fin-R1-7B
ships temperature $0.7$, top-$p$ $0.8$ and repetition penalty $1.05$. We depart from all three so
that every system is compared under identical decoding, and we do not claim the resulting numbers
equal what each vendor would obtain under its own defaults. Three API systems were additionally
run with an explicitly set, non-default reasoning budget, GPT-5.5 at \texttt{none},
Gemini-3.1-Pro at \texttt{low} and DeepSeek-V4-Pro at \texttt{high}, and these are marked in
Table~\ref{tab:leaderboard}. Items whose prompt length plus the completion budget would exceed
the context window are isolated as errors rather than silently truncated.

\subsection{Prompts and Structured Output Protocol}
\label{app:prompts}

Locally deployed models receive the prompt through their own chat template via
\texttt{apply\_chat\_template(..., tokenize=True, add\_generation\_prompt=True)}, and the
resulting token identifiers go directly to vLLM. This keeps each model aligned with its own
training format rather than imposing a single hand-written wrapper. The prompt is delivered as
one user-only message with no system message, which is the configuration under which the
finance-specialized checkpoints were tuned. API systems receive the identical user text through
their standard endpoints.

The prompt carries only the question content and its options. It never carries the keyed answer
or the provider rationale, both of which are present in the corpus and are used solely for
scoring and auditing. The prompt is byte-for-byte identical across all seventeen systems and
across all evaluation conditions, and is reproduced in Figure~\ref{fig:prompt-mcq}.

\begin{figure}[htbp]
\centering
\begin{promptbox}
\textbf{User}

\vspace{0.5em}
Solve this CFA/FRM multiple-choice finance question. Return JSON only with keys
\texttt{answer}, \texttt{rationale}, \texttt{confidence}. \texttt{answer} must be one of
A, B, C, D. Keep rationale concise.

\vspace{0.8em}
\textbf{Question:} \\
\texttt{\{content\}}

\vspace{0.5em}
\textbf{Options:} \\
\texttt{A. \{option\_A\}} \\
\texttt{B. \{option\_B\}} \\
\texttt{C. \{option\_C\}} \\
\texttt{D. \{option\_D\}}
\end{promptbox}
\caption{The single prompt used for every system and every condition. It is delivered as one
user-only message with no system message, through each local model's own chat template. CFA
items present three options and FRM items four, so the \texttt{D} line is emitted only when the
item has a fourth option. Option letters are taken from the item's own identifiers rather than
reassigned, so an item is never silently relabelled.}
\label{fig:prompt-mcq}
\end{figure}

\paragraph{Parsing.}
Responses are recovered with a cascade that stops at the first success. We first parse the whole
output as strict JSON. If that fails we retry under a permissive JSON reader that tolerates
trailing commas and unquoted keys. If that fails we extract an explicitly tagged answer field. If
no final message is present at all, which happens when a model emits only its reasoning channel,
the attempt is recorded as having no final. Anything still unrecovered is recorded as a parse
failure. For the two GPT-OSS models the Harmony analysis and final channels are separated at the
token level before parsing, so that reasoning text is never mistaken for an answer. Recovered
letters are uppercased before comparison.

Responses from which no option letter can be recovered are scored wrong rather than dropped, so
the denominator is the full benchmark for every system. This choice matters for the
finance-specialized group in particular, which produces the overwhelming majority of unparsable
responses, $1{,}224$ for Fin-R1-7B and $825$ for ODA-Fin-RL-8B against at most $231$ for any
general-purpose model. A protocol that discarded unparsable responses would report those systems
as substantially stronger than they are in use, since a caller receives nothing usable in those
cases. A length-truncated completion is likewise counted as a completed but incorrect attempt
rather than as an infrastructure error.

\paragraph{Monitoring and leakage control.}
Progress monitoring exposes only aggregate counters. No question text, option text, rationale,
reasoning channel or raw generation is written to any log, so that a monitoring artefact cannot
become an inadvertent release of the held-out half.
